\documentclass{article}

% Language setting
% Replace `english' with e.g. `spanish' to change the document language
\usepackage[english]{babel}

% Set page size and margins
% Replace `letterpaper' with `a4paper' for UK/EU standard size
\usepackage[letterpaper,top=2cm,bottom=2cm,left=3cm,right=3cm,marginparwidth=1.75cm]{geometry}

% Useful packages
\usepackage{amsmath}
\usepackage{graphicx}
\usepackage[colorlinks=true, allcolors=blue]{hyperref}

\title{Supplementary Material: Detailed Scoring Derivation}
\author{Or Forshmit \and Noam Kimhi \and Roee Tekoah \and Dor Stein \and Noam Korkos}
\date{}

\begin{document}
\maketitle
\noindent This document provides the full formal definition and derivation of the scoring scheme used in the main text.

\paragraph{Scoring rationale.}
The scoring framework is designed to detect edges and nodes that are anomalously similar relative to their species-pair background. 
Instead of relying on global similarity thresholds, we normalize similarities within each unordered species pair, thereby accounting for heterogeneity in evolutionary distance and sequence conservation across taxa. 
Robust statistics (median and MAD) are used to ensure stability under heavy-tailed similarity distributions and to suppress the influence of conserved proteins and outliers. This design avoids reliance on explicit phylogenetic reconstruction while still capturing deviations from expected cross-species similarity patterns.

\vspace{0.5em}
We use the exact implemented scoring equations. Let $e=(u,v)$ and let $p(e)$ be the unordered species pair of the proteins connected by edge $e$, i.e., the pair of species labels associated with the endpoints of $e$, ignoring order. All normalization is performed independently within each such species pair.
For each pair $p$ over edge set $E_p$ (default weight $w_e=$ Jaccard):
\begin{equation}
m_p=\mathrm{median}\{w_e:e\in E_p\},\quad
d_p=\mathrm{median}\{|w_e-m_p|:e\in E_p\},\quad
s_p=1.4826\,d_p+10^{-9}.
\end{equation}
Here $\mathrm{median}(\cdot)$ denotes the sample median. The quantity $d_p$ corresponds to the median absolute deviation (MAD), defined as $\mathrm{MAD}(X)=\mathrm{median}(|x - \mathrm{median}(X)|)$.
The scaling factor $1.4826$ ensures consistency with the standard deviation under a normal distribution, making $s_p$ comparable to a standard deviation estimate while retaining robustness to outliers. 
The small constant $10^{-9}$ prevents numerical instability when $d_p$ is very close to zero.

\vspace{0.5em}
Edge robust score:
\begin{equation}
z_e=
\begin{cases}
\dfrac{w_e-m_{p(e)}}{s_{p(e)}}, & |E_{p(e)}|\ge 30\\
\text{undefined}, & |E_{p(e)}|<30.
\end{cases}
\end{equation}
We use a global anomaly threshold $z_0=3.0$ to define high-anomaly edges. The requirement $|E_{p(e)}|\ge 30$ ensures sufficient sample size for stable estimation of the median and MAD; smaller samples yield unreliable baselines and are therefore excluded from anomaly scoring.
Thus, $z_e$ measures how many robust standard deviations the edge weight $w_e$ lies above the typical similarity for its species pair. 
Large positive values indicate unexpectedly high similarity relative to the background distribution, consistent with potential horizontal transfer signals.

Practically, if $z_e$ is undefined, the edge is not removed from the graph, but it is excluded from all $z$-based anomaly terms (node/component anomaly summaries and high-$z$ betweenness), so it does not contribute to anomaly-based statistics.
\vspace{0.5em}

Let $q_i$ denote the empirical fraction of nodes in component $C$ belonging to species $i$, and let $k_C$ be the number of distinct species in $C$. 
The normalized entropy $H_C^{\mathrm{norm}}$ is defined as the Shannon entropy of this discrete distribution divided by $\log k_C$, ensuring values in $[0,1]$.
The normalized neighbor-species entropy $H_u^{\mathrm{norm}}$ is defined analogously using the species distribution of neighbors of node $u$, normalized to lie in $[0,1]$.

For each component $C$:
\begin{equation}
H_C^{\mathrm{norm}}=
\begin{cases}
\dfrac{-\sum_i q_i\log q_i}{\log k_C}, & k_C\ge 2\\
0, & k_C<2,
\end{cases}
\qquad
\mathrm{high\_z\_frac}(C)=\frac{|\{e\in E_C:z_e\ge z_0\}|}{|E_C|}.
\end{equation}
\vspace{0.5em}

For each node $u$: degree $d(u)$; neighbor-species count $n_{\mathrm{sp}}(u)$; top-5 mean incident robust score $z_5(u)$, defined as the mean of the five largest defined $z_e$ (if fewer than five defined $z_e$ values exist, the mean is taken over available values) among edges incident to $u$; max neighbor-species fraction $f(u)$; normalized neighbor-species entropy $H_u^{\mathrm{norm}}$, clustering $cc(u)$, and betweenness $b(u)$ on the high-$z$ subgraph.

Hard gates are:
\begin{equation}
\begin{aligned}
&|V_{C_u}|\ge 10,\quad H_{C_u}^{\mathrm{norm}}\ge 0.75,\quad d(u)\ge 4,\quad n_{\mathrm{sp}}(u)\ge 2,\\
&z_5(u)\ \text{defined and}\ z_5(u)\ge 2.0,\\
&\big(n_{\mathrm{sp}}(u)\ge 3\ \text{or}\ b(u)>10^{-12}\big).
\end{aligned}
\end{equation}
Failure on any gate gives score $0$.
Threshold rationale (hard gates): $|V_{C_u}|\ge 10$ avoids unstable component statistics in very small components; $H_{C_u}^{\mathrm{norm}}\ge 0.75$ enforces strong taxonomic mixing; $d(u)\ge 4$ and $n_{\mathrm{sp}}(u)\ge 2$ suppress sparse single-link artifacts; $z_5(u)\ge 2.0$ requires clear robust anomaly signal; and $n_{\mathrm{sp}}(u)\ge 3$ or $b(u)>10^{-12}$ keeps either broadly connected multi-species nodes or bridge-like nodes.
\vspace{0.5em}

Within each component, robust baselines for $x\in\{z_5,f,d,b\}$ are
\begin{equation}
m_{x,C}=\mathrm{median}\{x(v):v\in V_C\},\quad
s_{x,C}=\max\!\left(1.4826\cdot \mathrm{MAD}\{x(v):v\in V_C\},\;10^{-3}(|m_{x,C}|+1)\right),
\end{equation}

The quantity $m_{x,C}$ is the median value of feature $x$ over nodes in component $C$, representing a typical baseline level within that component. 
The scale $s_{x,C}$ is a robust estimate of variability, defined using the MAD and lower-bounded by $10^{-3}(|m_{x,C}|+1)$ to avoid numerical instability in low-variance or degenerate cases.

Each feature is then standardized relative to its component-specific baseline, yielding a robust $z$-score:
If $Z_z(u)=Z_f(u)=Z_b(u)=0$, score is set to $0$.
\vspace{0.5em}


We define the following component-level and node-level factors. 
Here, $\rho(C)$ is component edge density; $P_{\mathrm{dens}}$ down-weights very dense components; $G_{\mathrm{ctx}}$ is the fraction of edges in $C$ with high anomaly score ($z\ge z_0$); $G_{\mathrm{size}}$ and $G_{\mathrm{deg}}$ are log-scaled boosts for larger component size and higher node degree; $P_{\mathrm{bridge}}$ rewards low local clustering (bridge-like topology); $P_{\mathrm{ent}}$ down-weights nodes with very low neighbor-species entropy; and $\mathrm{clip}_{[0,1]}(x)=\min(1,\max(0,x))$ truncates values to the interval $[0,1]$.
\begin{equation}
\rho(C)=\frac{2|E_C|}{|V_C|(|V_C|-1)},\;
P_{\mathrm{dens}}=\max(10^{-6},(1-\mathrm{clip}_{[0,1]}(\rho))^{2}),\;
G_{\mathrm{ctx}}=\max(\mathrm{high\_z\_frac},10^{-6}),
\end{equation}
\begin{equation}
G_{\mathrm{size}}=\log(1+|V_C|),\;
G_{\mathrm{deg}}=\log(1+d(u)),\;
P_{\mathrm{bridge}}=\max(10^{-6},(1-\mathrm{clip}_{[0,1]}(cc(u)))),
\end{equation}
\begin{equation}
P_{\mathrm{ent}}=
\begin{cases}
0.1+0.9\cdot \dfrac{H_u^{\mathrm{norm}}}{0.15}, & H_u^{\mathrm{norm}}<0.15\\
1, & H_u^{\mathrm{norm}}\ge 0.15.
\end{cases}
\end{equation}

\vspace{0.5em}
\textbf{Final score.} Let $Z_z,Z_f,Z_d,Z_b$ denote the positive-capped robust standardized node features corresponding to $(z_5,f,d,b)$, respectively. The score is constructed in log-space to combine multiplicative factors additively, improving numerical stability and interpretability. Each term contributes independently, enabling explicit control over the influence of node-level anomaly, structural context, and component-level properties.
The density penalty $P_{\mathrm{dens}}$ down-weights highly dense (near-clique) components, which are more likely to arise from conserved protein families rather than localized transfer events.
The context term $G_{\mathrm{ctx}}$ reflects the fraction of high-anomaly edges in the component, favoring components enriched with unusually strong cross-species similarities.
\begin{equation}
\begin{aligned}
\log S(u)={}&\log G_{\mathrm{size}}+\log G_{\mathrm{ctx}}+\log G_{\mathrm{deg}}
+\log P_{\mathrm{dens}}+\log P_{\mathrm{bridge}}+\log P_{\mathrm{ent}}\\
&+1.0\log(1+Z_z)+2.0\log(1+Z_f)+0.5\log(1+0.25Z_d)+1.0\log(1+Z_b),
\end{aligned}
\end{equation}
with $S(u)=\exp(\log S(u))$. The coefficients on $\log(1+Z_x)$ terms control the relative importance of each feature: anomaly strength ($Z_z$), cross-species connectivity ($Z_f$), degree ($Z_d$), and structural centrality ($Z_b$). 
These weights are fixed hyperparameters reflecting empirical prioritization of signals indicative of horizontal transfer.
\vspace{0.5em}

\paragraph{Additional analyzed features beyond direct score terms.}
Beyond the direct multiplicative terms in $S(u)$, the pipeline also computes and reports additional descriptive features for interpretation and diagnostics. Edge-level outputs include raw Jaccard, shared-$k$-mer count, species-pair baseline fields (pair median, pair MAD, pair sample size), and robust $z$. Node-level outputs include top-1/top-5 Jaccard summaries, top-1/top-5 shared-$k$-mer summaries, max incident $z$, effective species count $n_{\mathrm{eff}}=\exp(H)$, and participation coefficient. Component-level outputs include max component $z$ and top-10 mean component $z$. These are written to result tables and used in reporting/visualization even when not direct multiplicative factors in $S(u)$.

All thresholds above are fixed hyperparameters in code (not learned and not estimated per component). We keep a global $z_0=3.0$ for comparability across components, require at least 30 edges per species pair to stabilize MAD-based normalization, and use conservative node/component gates selected during exploratory calibration to maintain high selectivity with a non-zero candidate set.
\subsubsection*{Hyperparameters}

\textbf{Candidate generation:}
$k=6$, min length $=50$, max posting list $=100$, min shared $k$-mers $=6$, top matches $=10$, cross-species only.

\textbf{Pre-pruning:}
species-pair quantile $q=0.9$, degree cap $=20$, min Jaccard $=0.1$, taxonomy filter $[2,4]$.

\textbf{Normalization:}
min edges per species pair $=30$, scale $1.4826\cdot \mathrm{MAD} + 10^{-9}$, anomaly threshold $z_0=3$.

\textbf{Gating:}
$|V_C|\ge10$, $H_C^{\mathrm{norm}}\ge0.75$, $d(u)\ge4$, $n_{\mathrm{sp}}(u)\ge2$, $z_5(u)\ge2$, plus $(n_{\mathrm{sp}}\ge3$ or $b>10^{-12})$.

\textbf{Scoring weights:}
$w_z=1$, $w_f=2$, $w_d=0.5$, $w_b=1$, $z_{\max}=10$.

\end{document}
